% Converted from README.md
\subsection{C++ - QtTest, gcov, lcov}\label{c---qttest-gcov-lcov}

\subsubsection{Pisanje testova jedinica koda pomoću QtTest radnog okvira}\label{pisanje-testova-jedinica-koda-pomoux107u-qttest-radnog-okvira}

QtCreator kao razvojno okruženje pruža mogućnost pisanja jediničnih testova. To je moguće učiniti pisanjem projekta tipa \texttt{Qt\ Unit\ Test} na više načina, u zavisnosti od toga kako želimo da organizujemo stvarni i test projekat.

\paragraph{Uključivanjem biblioteke}\label{ukljuux10divanjem-biblioteke}

Prvi način kako možemo kreirati test projekat koji testira već postojeći projekat je tako što kreiramo novi projekat tipa \texttt{Qt\ Unit\ Test}:

\begin{minted}{text}
New Project -> Other Project -> Qt Unit Test
\end{minted}

Pošto želimo da pišemo testove za već postojeći projekat, treba da dodamo lokaciju projekta opcijom \texttt{Add\ Existing\ Directory}.

Unutar test projekta će se kreirati klasa koja nasleduje \texttt{QObject}. Ispod modifikatora \texttt{private\ Q\_SLOTS} pišemo bar jednu test funkciju. Trebalo bi da su nam već ponuđene funkcije npr. \texttt{testCase1()}.

Dodatno, postoje i četiri privatne metode koje se ne tretiraju kao test funkcije, ali će ih test radni okvir izvršavati bilo kada inicijalizuje ili čisti za celim testom ili trenutnom test funkcijom:

\begin{itemize}
\tightlist
\item
  \texttt{initTestCase()} će biti pozvana pre izvršavanja prve test funkcije
\item
  \texttt{cleanupTestCase()} će biti pozvana nakon izvršavanja poslednje test funkcije
\item
  \texttt{init()} će biti pozvana pre svakog poziva test funkcije
\item
  \texttt{cleanup()} će biti pozvana nakon svakog izvršavanja test funkcije
\end{itemize}

Ukoliko se \texttt{initTestCase()} ne izvrši uspešno, nijedna test funkcija neće biti izvršavana. Ako \texttt{init()} funkcija ne prođe, njena prateća test funkcija se neće izvršiti, ali će se nastaviti sa sledećom.

\paragraph{Qt Subdirs projekat}\label{qt-subdirs-projekat}

Drugi način da testiramo projekat je da kreiramo nov \texttt{Qt\ Subdirs} projekat u koji ćemo uključiti postojeći i novi \texttt{Qt\ Unit\ Test} projekat. Ovaj način organizacije je koristan ukoliko pišemo projekat od nule i želimo da usput pišemo i testove.

\subsubsection{Pisanje testova}\label{pisanje-testova}

QtTest radni okvir pruža makroe za testiranje, i neke od njih ćemo koristiti da testiramo naš kalkulator:

\begin{itemize}
\tightlist
\item
  \texttt{QCOMPARE}
\item
  \texttt{QVERIFY}
\item
  \texttt{QVERIFY\_EXCEPTION\_THROWN}
\end{itemize}

\subsubsection{Analiza pokrivenosti}\label{analiza-pokrivenosti}

Uz \texttt{gcc} kompajler dolazi i \texttt{gcov} alat za određivanje pokrivenosti koda prilikom izvršavanja programa (engl. \emph{code coverage}). Koristi se zajedno sa \texttt{gcc} kompajlerom da bi se analizirao program i utvrdilo kako se može kreirati efikasniji, brži kod i da bi se testovima pokrili delovi programa.

Alat \texttt{gcov} se može koristiti kao alat za profajliranje u cilju otkrivanja dela koda čija bi optimizacija najviše doprinela efikasnosti programa. Korišćenjem \texttt{gcov}-a možemo saznati koje su naredbe, linije, grane, funkcije itd. izvršene i koliko puta. Zarad lepše reprezentacije rezultata detekcije pokrivenosti koda izvršavanjem test primera, koristimo alat \texttt{lcov}.

Prilikom kompilacije neophodno je koristiti dodatne opcije kompajlera koje omogućavaju snimanje koliko je puta koja linija, grana i funkcija izvršena. Ti podaci se čuvaju u datotekama ekstenzije \texttt{.gcno} za svaku datoteku sa izvornim kodom. One će kasnije biti korišćene za kreiranje izveštaja o pokrivenosti koda.

\begin{minted}{bash}
$ g++ -g -Wall -fprofile-arcs -ftest-coverage -O0 main.cpp -o test
\end{minted}

Alternativno:

\begin{minted}{bash}
g++ -g -Wall --coverage -O0 main.cpp -o test
\end{minted}

Nakon izvršavanja test programa, informacije o pokrivenosti biće sačuvane u datoteci tipa \texttt{.gcda}, ponovo za svaku datoteku sa izvornim kodom. Pokrenimo alat \texttt{lcov} da bismo dobili čitljiviji prikaz rezultata:

\begin{minted}{bash}
$ lcov --rc lcov_branch_coverage=1 -c -d . -o coverage-test.info
\end{minted}

Opcija:

\begin{itemize}
\tightlist
\item
  \texttt{-\/-rc\ lcov\_branch\_coverage=1} uključuje određivanje pokrivenosti grana, koje podrazumevano nije uključeno
\item
  \texttt{-c} kreiranje pokrivenosti
\item
  \texttt{-d\ .} koristi tekući direktorijum, jer u našem slučaju on sadrži potrebne \texttt{.gcda} i \texttt{.gcno} datoteke
\item
  \texttt{-o\ coverage-test.info} zadaje naziv izlazne datoteke sa izveštajem koji treba da ima ekstenziju \texttt{.info}
\end{itemize}

Možemo neke datoteke isključiti iz analize pokrivenosti. Na primer, biblioteke jezika koje ne testiramo mogu nam samo zamagliti pokrivenost koja nas zanima - pokrivenost funkcionalnosti koje testiramo.

\begin{minted}{bash}
$ lcov --rc lcov_branch_coverage=1 \
    -r coverage.info ’/usr/*’ ’/opt/*’ ’*.moc’ \
    -o coverage-filtered.info
\end{minted}

Opcija \texttt{-r\ coverage.info} uklanja iz prethodno dobijenog izveštaja \texttt{coverage.info} izvorne datoteke koje odgovaraju nekom od šablona koji su navedeni kao argumenti opcije.

Alat \texttt{lcov} ima podalat \texttt{genhtml} koji na osnovu prethodno generisanog izveštaja pravi \texttt{.html} datoteke za jednostavniji pregled. Potrebno je izvršiti naredbu:

\begin{minted}{bash}
$ genhtml --rc lcov_branch_coverage=1 -o Reports coverage-filtered.info
\end{minted}

Opcija \texttt{-o\ Reports} određuje naziv direktorijuma koji će biti kreiran i popunjen generisanim \texttt{.html} dokumentima.

Izveštaj možemo otvoriti u Web pretraživaču, npr.

\begin{minted}{bash}
$ firefox Reports/index.html
\end{minted}

Ukoliko ne obrišemo \texttt{.gcda} datoteke od prethodnih pokretanja programa, prikaz pokrivenosti će uključiti sve, zbirno.

\subsubsection{Kreiranje izveštaja uz pomoć skripte}\label{kreiranje-izveux161taja-uz-pomoux107-skripte}

Postupak kreiranja izveštaja nakon kompilacije programa može se automatizovati pokretanjem bash skripte \texttt{generateCodeCoverageReport.sh}:

\begin{minted}{bash}
$ ./generateCodeCoverageReport.sh . test data
\end{minted}

Argumenti:

\begin{itemize}
\tightlist
\item
  \texttt{.} - Direktorijum u kom se nalaze potrebne \texttt{.gcda} i \texttt{.gcno} datoteke i izvršni program. U našem slučaju to je tekući direktorijum. Inače bismo navodili relativnu ili apsolutnu putanju do potrebnog direktorijuma.
\item
  \texttt{test} - Drugi argument treba da je naziv izvršne verzije programa koja će se pokretati.
\item
  \texttt{data} - Naziv direktorijuma u koji će alati \texttt{lcov} i \texttt{genhtml\ }upisivati svoje rezultate. Ukoliko se ne navede, sve će se upisivati u tekući direktorijum skripte. Ukoliko navedeni direktorijum ne postoji, biće kreiran.
\end{itemize}

Skript briše \texttt{.gcda} datoteke od prethodnih pokretanja programa. Prikazana pokrivenost je samo za poslednje pokretanje programa.

\subsubsection{Integracija pokrivenosti u QtCreator}\label{integracija-pokrivenosti-u-qtcreator}

Ukoliko želimo da kreiramo pokrivenost koda prilikom izvršavanja postojećeg Qt projekta, potrebno je da definiciji projekta (datoteci tipa \texttt{.pro}) dopišemo:

\begin{minted}{bash}
QMAKE_CXXFLAGS += -g -Wall -fprofile-arcs -ftest-coverage -O0
QMAKE_LFLAGS += -g -Wall -fprofile-arcs -ftest-coverage -O0

LIBS += \
    -lgcov
\end{minted}

Opcija kompajlera \texttt{-\/-coverage} menja opcije \texttt{-fprofile-arcs\ -ftest-coverage} i linkovanje sa \texttt{-lgcov}. Nakon pokretanja projekta, \texttt{.gcda} i \texttt{.gcno} datoteke i izvršivi program biće u direktorijumu gde se nalaze ostali artefakti prevođenja (podrazumevano u direktorijumu sa prefiksom \texttt{build\_}). Izveštaj možemo potom napraviti prema ranije prikazanom postupku.
